
i.e. qui rend @@M0@@ nul. Comme la formation de @@M1@@ commute au changement de base, on voit qu'il est représenté par l'ouvert @@M2@@ de @@M3@@. On peut supposer donc que @@M4@@ est déjà un épimorphisme. Mais utilisant maintenant l'hypothèse que @@M5@@ est @@M6@@-plat, on trouve que la formation de @@M7@@ commute au changement de base, et par suite le foncteur envisagé est représenté par l'ouvert @@M8@@ de @@M9@@.



% Expose XIII local macros
\providecommand{\Ass}{\operatorname{Ass}}
\providecommand{\Hom}{\operatorname{Hom}}
\providecommand{\Ext}{\operatorname{Ext}}
\providecommand{\Proj}{\operatorname{Proj}}
\providecommand{\NS}{\operatorname{NS}}
\providecommand{\Tors}{\operatorname{Tors}}
\providecommand{\Picnum}{\operatorname{Pic}^{\mathrm{num}}}
\providecommand{\Piczero}{\operatorname{Pic}^{0}}
\providecommand{\Pictau}{\operatorname{Pic}^{\tau}}
\providecommand{\length}{\operatorname{length}}
\providecommand{\etale}{\text{\'etale}}


